\chapter{Zadatak 1a --- je li problem konveksan}
\label{ch:z1-konveksnost}

\section{Što se ovdje traži}

Treba dokazati --- ne samo tvrditi --- da problem pripada klasi konveksnih
optimizacijskih problema. To znači pokazati dvije stvari odvojeno: da je funkcija
cilja konveksna i da je dopustivi skup konveksan.

\begin{zadatak}{Zadatak 1, podzadatak a)}
a) Je li ovo konveksan optimizacijski problem? Obrazložiti odgovor.
\end{zadatak}

\begin{predznanje}
\begin{itemize}[nosep,leftmargin=*]
  \item Što je konveksan skup i konveksna funkcija --- \S\ref{sec:prostori}, Poglavlje~\ref{ch:rjecnik}.
  \item Razlika između $\leq$, $\psd$ i $\pd$ --- \S\ref{sec:relacije}.
  \item Računanje svojstvenih vrijednosti $2\times2$ matrice --- \S\ref{sec:svojstvene}.
  \item Provjera definitnosti preko minora --- \S\ref{sec:definitnost}.
  \item Hessian kvadratne funkcije --- \S\ref{sec:gradijent-hessian}.
\end{itemize}
\end{predznanje}

\section{Definicija koju treba zadovoljiti}
\label{sec:z1a-definicija}

\begin{izvor}{Predavanje 2 --- Definicije, klasifikacija, konveksnost, slajd 70/71 (PDF str.~71)}
\slajd[0.86]{p02-konveksan-problem}

\textbf{Transkripcija.} \emph{Konveksan optimizacijski problem}
\[
\begin{aligned}
\text{minimiziraj} \quad & f(x) \\
\text{uz ograničenja} \quad & h_i(x) = 0, && i = 1, \dots, m \\
& g_j(x) \leq 0, && j = 1, \dots, m
\end{aligned}
\]
gdje su $f(x)$ i $g_j(x)$, $j = 1, \dots, m$, konveksne funkcije, a $h_i(x)$,
$i = 1, \dots, m$, afine funkcije.

\textbf{Na što treba paziti u ovoj definiciji:} uvjet nije simetričan. Od
funkcija nejednakosti $g_j$ traži se samo konveksnost, ali od funkcija jednakosti
$h_i$ traži se \emph{afinost}, što je jače. Razlog: skup $\{x \mid h(x) = 0\}$
konveksne funkcije $h$ općenito nije konveksan (ploha, ne poluprostor), dok
$\{x \mid g(x) \leq 0\}$ konveksne $g$ uvijek jest.
\end{izvor}

\begin{pojam}{Konveksan optimizacijski problem --- kontrolna lista}
Problem je konveksan ako vrijedi \emph{sve troje}:
\begin{enumerate}[nosep,leftmargin=*]
  \item funkcija cilja $f$ je konveksna;
  \item svaka funkcija ograničenja nejednakosti $g_j$ je konveksna;
  \item svaka funkcija ograničenja jednakosti $h_i$ je afina.
\end{enumerate}
U Zadatku 1 ograničenja jednakosti nema, pa se treća točka ispušta.
\end{pojam}

\section{Konveksan skup --- definicija na kojoj sve stoji}
\label{sec:z1a-konveksan-skup}

\begin{izvor}{Predavanje 2 --- Definicije, klasifikacija, konveksnost, slajd 55/71 (PDF str.~56)}
\slajd[0.86]{p02-konveksan-skup}

\textbf{Transkripcija.} Konveksan skup $S$ definiran je sljedećim svojstvom:
dužina između bilo koje dvije različite točke iz $S$ se i sáma u potpunosti
nalazi u $S$.

Dužina između dvije točke $x_1, x_2 \in \R^n$, $x_1 \neq x_2$, je skup $D$, gdje je
\[
D \coloneqq \{\, x \mid x = \alpha x_1 + (1-\alpha)x_2 \,\}, \qquad 0 \leq \alpha \leq 1 .
\]
\emph{Konveksan skup:} skup $S$ je konveksan skup ako
\[
x_1, x_2 \in S \;\Longrightarrow\; \alpha x_1 + (1-\alpha)x_2 \in S
\quad \text{za sve } \alpha \in [0,1].
\]

\textbf{Kako se to čita naglas:} ,,alfa iks jedan plus jedan minus alfa iks dva``
je opća točka na spojnici dviju točaka; za $\alpha = 1$ dobiva se $x_1$, za
$\alpha = 0$ dobiva se $x_2$, a za međuvrijednosti sve točke između. Uvjet dakle
kaže: uzmi bilo koje dvije točke skupa, spoji ih ravnom crtom --- crta ne smije
nigdje izaći iz skupa.

\textbf{Zašto je to ovdje bitno:} poluprostor očito zadovoljava taj uvjet, jer
ako su obje krajnje točke s iste strane pravca, cijela je spojnica s te strane.
Odatle konveksnost svakog $\feas_i$, a onda i presjeka.
\end{izvor}

\section{Kriterij drugog reda}
\label{sec:z1a-kriterij}

\begin{izvor}{Predavanje 2 --- Definicije, klasifikacija, konveksnost, slajd 69/71 (PDF str.~70)}
\slajd[0.86]{p02-uvjeti-drugog-reda}

\textbf{Transkripcija.} \emph{Konveksne funkcije --- uvjeti drugog reda.} Neka je
$f: \R^n \to \R$ diferencijabilna funkcija s konveksnom domenom. Tada vrijedi:
\begin{itemize}[nosep,leftmargin=*]
  \item $f$ je konveksna funkcija \textbf{ako i samo ako} $\grad^2 f(x) \psd 0$
        za sve $x$ u domeni;
  \item ako je $\grad^2 f(x) \pd 0$ za sve $x$ u domeni, tada je $f$
        \textbf{striktno} konveksna funkcija.
\end{itemize}

\textbf{Pažljivo pročitaj razliku:} prva tvrdnja je ekvivalencija (,,ako i samo
ako``), druga je samo implikacija u jednom smjeru. Striktno konveksna funkcija ne
mora imati svugdje pozitivno definitan Hessian --- primjer je $f(x) = x^4$, koja
je striktno konveksna, a $f''(0) = 0$. Zato se iz $\grad^2 f \psd 0$ smije
zaključiti konveksnost, ali iz striktne konveksnosti ne smije se zaključiti
$\grad^2 f \pd 0$.
\end{izvor}

\section{Izvod za naš problem}
\label{sec:z1a-izvod}

\begin{korak}[Korak 1: Hesseova matrica funkcije cilja]
Funkcija cilja je kvadratna, $f(x) = \tfrac12 x\trans Q x + c\trans x$, pa je po
\eqref{eq:grad-hess-kvadratne}
\begin{equation}
\grad^2 f(x) = Q = \begin{bmatrix} 2 & -0{,}3 \\ -0{,}3 & 4 \end{bmatrix}
\qquad \text{za svaki } x \in \R^2 .
\end{equation}
Hessian je konstantan, pa je dovoljno ispitati ga jednom --- nema potrebe za
provjerom ,,u svakoj točki``, kako doslovno traži kriterij.
\end{korak}

\begin{korak}[Korak 2: svojstvene vrijednosti]
Karakteristična jednadžba \eqref{eq:karakteristicna-2x2} uz
$\tr Q = 6$ i $\det Q = 2\cdot4 - (-0{,}3)^2 = 8 - 0{,}09 = 7{,}91$:
\begin{equation}
\det(\lambda I - Q) = (\lambda - 2)(\lambda - 4) - 0{,}09
= \lambda^2 - 6\lambda + 7{,}91 = 0,
\end{equation}
\begin{equation}
\lambda_{1,2} = \frac{6 \pm \sqrt{36 - 4\cdot 7{,}91}}{2}
             = \frac{6 \pm \sqrt{4{,}36}}{2}
\;\Longrightarrow\;
\lambda_1 = 1{,}955969, \quad \lambda_2 = 4{,}044031 .
\end{equation}
Obje su strogo pozitivne, dakle $Q \pd 0$.
\end{korak}

\begin{korak}[Korak 3: neovisna provjera preko minora]
Po kriteriju s Predavanja 2 (\S\ref{sec:definitnost}):
\begin{equation}
D_1 = Q_{11} = 2 > 0, \qquad
D_2 = \det Q = 7{,}91 > 0 .
\end{equation}
Oba su pozitivna, pa je $Q \pd 0$ --- isti zaključak, dobiven neovisnim putem.
\end{korak}

\begin{korak}[Korak 4: konveksnost funkcije cilja]
Iz $\grad^2 f(x) = Q \pd 0$ za sve $x$ i kriterija drugog reda slijedi da je $f$
\textbf{striktno} konveksna na cijelom $\R^2$.
\end{korak}

\begin{korak}[Korak 5: konveksnost ograničenja]
Funkcije $g_1(x) = -x_1 - x_2 + 10$ i $g_2(x) = x_1 - x_2 - 3$ su \emph{afine}:
oblika $a\trans x + \beta$. Afina funkcija je konveksna (i konkavna istovremeno),
jer joj je Hessian nul-matrica, a $0 \psd 0$. Dakle oba ograničenja nejednakosti
zadovoljavaju uvjet iz definicije.
\end{korak}

\begin{korak}[Korak 6: konveksnost dopustivog skupa]
Svako pojedinačno ograničenje definira poluprostor
$\feas_i = \{x \mid g_i(x) \leq 0\}$, a poluprostor je konveksan skup. Po svojstvu
s Predavanja 2 (\S\ref{sec:poluprostor}, slajd 58/71) presjek konveksnih skupova
je konveksan, pa je
\begin{equation}
\feas = \feas_1 \cap \feas_2
\end{equation}
zatvoren konveksan skup.
\end{korak}

\begin{korak}[Korak 7: zaključak i još dvije tvrdnje koje se očekuju]
Funkcija cilja je konveksna, ograničenja su konveksna, ograničenja jednakosti
nema --- \textbf{problem je konveksan}. Uz to:
\begin{itemize}[leftmargin=*]
  \item \textbf{Egzistencija minimuma.} Kako je $Q \pd 0$, funkcija $f$ je
        \emph{koercivna}: $f(x) \to +\infty$ kad $\norm{x} \to \infty$. Skup
        $\feas$ je zatvoren i neprazan (npr.\ $x = (7,6) \in \feas$), pa minimum
        postoji, iako je $\feas$ neograničen.
  \item \textbf{Jedinstvenost.} Stroga konveksnost daje jedinstven minimizator:
        dvije različite optimalne točke dale bi na spojnici manju vrijednost od
        obiju, što je proturječje.
\end{itemize}
\end{korak}

\section{Slaterov uvjet}
\label{sec:z1a-slater}

Ovo zadatak ne traži izrijekom, ali bez toga se pod (b) ne smije reći da je
KKT točka optimum.

\begin{izvor}{Predavanje 4 --- Lagrangeova dualnost, slajd 11/26 (PDF str.~16)}
\slajd[0.86]{p04-jaka-dualnost}

\textbf{Transkripcija.} \emph{Teorem Lagrangeove dualnosti.}
\begin{itemize}[nosep,leftmargin=*]
  \item Slaba dualnost uvijek vrijedi: $\dopt \leq \popt$.
  \item Kad je primaran problem konveksan i kad su zadovoljeni Slaterovi uvjeti,
        vrijedi jaka dualnost: $\dopt = \popt$.
\end{itemize}
Jaka dualnost u kompaktnoj formi:
\[
\max_{\{\lambda,\mu \mid \mu \geq 0\}} \ \inf_x \bigl(f(x) + \lambda\trans h(x) + \mu\trans g(x)\bigr)
= \inf_{\{x \mid g(x) \leq 0,\, h(x) = 0\}} f(x).
\]

\textbf{Zašto nam ovo treba već sada:} jaka dualnost je pretpostavka pod kojom
KKT uvjeti postaju \emph{dovoljni}, a ne samo nužni. Bez nje bi rješenje KKT
sustava bilo tek kandidat.
\end{izvor}

\begin{pojam}{Slaterov uvjet}
Za konveksan problem Slaterov uvjet je zadovoljen ako postoji \emph{strogo}
dopustiva točka, dakle $\bar{x}$ takva da je $g_j(\bar{x}) < 0$ za sva
ograničenja nejednakosti (stroga nejednakost, ne $\leq$).
\end{pojam}

Za naš problem uzmimo $\bar{x} = (7,\, 6)$:
\begin{equation}
g_1(\bar{x}) = -7 - 6 + 10 = -3 < 0, \qquad
g_2(\bar{x}) = 7 - 6 - 3 = -2 < 0 .
\end{equation}
Slaterov uvjet je zadovoljen. Dakle vrijedi jaka dualnost, a KKT uvjeti su
\textbf{nužni i dovoljni} za optimalnost.

\section{Kako je to riješeno u seminaru}
\label{sec:z1a-seminar}

\begin{seminar}{Seminar v1.2, §2.2 --- ,,(a) Je li problem konveksan?``}
Hesseova matrica funkcije cilja jednaka je matrici $Q$. Njezine svojstvene
vrijednosti slijede iz karakteristične jednadžbe
\[
\det(\lambda I - Q) = (\lambda - 2)(\lambda - 4) - 0.09 = \lambda^2 - 6\lambda + 7.91 = 0,
\]
\[
\lambda_{1,2} = \frac{6 \pm \sqrt{4.36}}{2} \implies \lambda_1 = 1.955969, \quad \lambda_2 = 4.044031.
\]
Obje su strogo pozitivne, pa je $Q \succ 0$ i funkcija cilja je strogo konveksna.
Funkcije ograničenja $g_1$ i $g_2$ su afine, dakle konveksne, pa je dopustivi skup
\[
\mathcal{F} = \{x \in \mathbb{R}^2 \mid Gx - h \leq 0\}
\]
presjek dvaju zatvorenih poluprostora, tj.\ zatvoren konveksan skup.

Riječ je, dakle, o konveksnom optimizacijskom problemu. Kako je $Q \succ 0$,
funkcija $f$ je koercivna ($f(x) \to \infty$ kada $\|x\| \to \infty$), a skup
$\mathcal{F}$ je zatvoren i neprazan, pa minimum postoji; stroga konveksnost daje
njegovu jedinstvenost. Za točku $x = [7, 6]^{\mathsf{T}}$ vrijedi $g_1(x) = -3 < 0$
i $g_2(x) = -2 < 0$, pa je Slaterov uvjet zadovoljen. Prema tome su KKT uvjeti
nužni i dovoljni za optimalnost i vrijedi jaka dualnost.
\end{seminar}

Numerička potvrda iz MATLAB-a (\texttt{docs/code/\allowbreak task1\_output.txt}):

\begin{lstlisting}[style=outputstyle]
--- (a) Konveksnost ---
eig(Q) = 1.955969, 4.044031
det(Q) = 7.910000,  Q > 0 : 1
\end{lstlisting}

\section{Isti podzadatak na vježbama}
\label{sec:z1a-vjezbe}

Na Vježbama 3 profesor je riješio zadatak istoga tipa, s podzadacima a)--f)
raspoređenima točno kao u ovom seminaru. Vrijedi vidjeti koliko je kratko
očekivano obrazloženje pod a).

\begin{izvor}{Vježbe 3 --- Primjer 2: Dualnost, KKT, osjetljivost, slajd 17/36 (PDF str.~22)}
\slajd[0.86]{v3-primjer2-konveksnost}

\textbf{Transkripcija.} \emph{a) Da li je ovo konveksan opt.\ problem?}
Funkcije ograničenja nejednakosti su konveksne (ovdje su one afine, pa stoga i
konveksne). Hessian:
\[
H = \begin{bmatrix}
\dfrac{\partial^2 f}{\partial x_1^2} & \dfrac{\partial^2 f}{\partial x_1 \partial x_2} \\[8pt]
\dfrac{\partial^2 f}{\partial x_2 \partial x_1} & \dfrac{\partial^2 f}{\partial x_2^2}
\end{bmatrix}
= \begin{bmatrix} 2 & 0 \\ 0 & 2 \end{bmatrix} .
\]

\textbf{Usporedba s našim zadatkom:} u profesorovu primjeru funkcija cilja je
$(x_1-3)^2 + (x_2-2)^2$, pa je Hessian dijagonalan i definitnost se vidi
napamet. Kod nas mješoviti član $-0{,}3\,x_1x_2$ daje vandijagonalne elemente, pa
se definitnost mora stvarno izračunati --- ali struktura obrazloženja je
identična: prvo ograničenja, pa Hessian, pa zaključak.
\end{izvor}

\section{Zamke}
\label{sec:z1a-zamke}

\begin{zamka}[Zamka 1: ,,funkcija cilja je konveksna`` nije cijeli odgovor]
Zadatak pita je li \emph{problem} konveksan, ne funkcija. Odgovor mora pokriti i
dopustivi skup. Student koji stane na $Q \pd 0$ odgovorio je na pola pitanja.
\end{zamka}

\begin{zamka}[Zamka 2: pozivanje na predznake elemenata matrice]
$Q$ ima negativan element $-0{,}3$. To ne govori ništa o definitnosti
(\S\ref{sec:relacije}). Jedini valjani argumenti su svojstvene vrijednosti,
vodeći glavni minori ili izravno $z\trans Q z > 0$.
\end{zamka}

\begin{zamka}[Zamka 3: prešućena egzistencija minimuma]
Dopustivi skup je \emph{neograničen}. Konveksnost sama po sebi ne jamči da se
minimum postiže --- primjer: $\min e^{x}$ na $\R$ je konveksan problem bez
minimuma. Ovdje minimum postoji jer je $f$ koercivna, a to treba reći.
\end{zamka}

\begin{zamka}[Zamka 4: brkanje ,,konveksna`` i ,,striktno konveksna``]
Iz $Q \pd 0$ slijedi \emph{striktna} konveksnost, a odatle jedinstvenost
rješenja. Da je bilo samo $Q \psd 0$, problem bi i dalje bio konveksan, ali bi
optimalnih točaka moglo biti beskonačno mnogo. Razlika je bitna, jer se u (b)
tvrdi da je nađeno rješenje \emph{to} rješenje.
\end{zamka}

\section{Modelski odgovor}
\label{sec:z1a-odgovor}

\begin{izgovor}[Odgovor u pet rečenica]
,,Da, problem je konveksan, i to iz dva odvojena razloga.

Prvo, funkcija cilja: zapisao sam je kao jedna polovina iks transponirano ku iks
plus ce transponirano iks, gdje je ku dva, minus nula tri, minus nula tri,
četiri. Hesseova matrica kvadratne funkcije jednaka je matrici ku i konstantna
je. Njezine svojstvene vrijednosti su jedan zarez devet pet šest i četiri zarez
nula četiri, obje strogo pozitivne, a isto potvrđuju i vodeći glavni minori, dva
i sedam zarez devet jedan. Dakle ku je pozitivno definitna i funkcija cilja je
striktno konveksna.

Drugo, dopustivi skup: oba ograničenja su afina, dakle konveksna, a svako
definira poluprostor. Presjek konveksnih skupova je konveksan, pa je dopustivi
skup zatvoren konveksan skup. Ograničenja jednakosti nema.

Dodao bih još dvoje. Minimum postoji jer je funkcija koercivna, a dopustivi skup
zatvoren i neprazan --- to je važno jer je skup neograničen. I jedinstven je,
zbog stroge konveksnosti.

Napokon, Slaterov uvjet je zadovoljen: točka sedam, šest daje $g_1 = -3$ i
$g_2 = -2$, oboje strogo negativno. Zato vrijedi jaka dualnost i KKT uvjeti su
nužni i dovoljni, što koristim u sljedećem podzadatku.``
\end{izgovor}
